A finals de l'any 2025, la Generalitat de Catalunya va llançar a l'espai un nanosatèl·lit que descriu una òrbita circular al voltant de la Terra a una altura aproximada de 500~km per sobre de la superfície terrestre.

\begin{apartat}{1,25}
A partir de la llei de la gravitació universal, deduïu l'expressió del període orbital en funció del radi orbital i de la massa de la Terra, i calculeu el període orbital del nanosatèl·lit.
\end{apartat}

\begin{apartat}{1,25}
Sabent que la massa del nanosatèl·lit és d'11,1~kg, calculeu l'increment d'energia potencial gravitatòria del satèl·lit en passar de la superfície terrestre a l'òrbita. Calculeu també l'energia mecànica del satèl·lit a l'òrbita. Es tracta d'una òrbita tancada? Justifiqueu la resposta.
\end{apartat}

\dades{%
  $G = 6,67\times 10^{-11}\ \mathrm{N\,m^2\,kg^{-2}}$.\\
  Massa de la Terra: $M_\mathrm{T} = 5,98\times 10^{24}\ \mathrm{kg}$.\\
  Radi de la Terra: $R_\mathrm{T} = 6\,370\ \mathrm{km}$.}
